\subsection{Experimental setup}
\label{sec:experiments}

We report partial future-prediction and downstream results in
\cref{tab:future_prediction,tab:downstream_tasks}. MIMIC-CXR~\citep{johnson2019mimiccxr}
supplies images and reports, with source-available MIMIC-IV history for forecasting.
Forecast training/validation/test sets contain 16,000/230/297 pairs; the test set
contains 94 patients. Adapted representation models use 2,400 updates with
batch size 32. The three legacy Qwen3.5-9B forecast conditions complete the same budget
from a shared 9B Stage~1 initialization. Dense tasks use 4,096/249/447 images and
20 training epochs; human lung Dice is evaluated separately on 138 Montgomery
images. The tables distinguish zero-shot checkpoints, frozen-feature task heads,
and trained forecast adaptations. The completed 9B slot adaptation uses final-layer
language queries; it does not implement the proposed four fusion-layer plus four
native-vision slots. The maintained unified implementation uses Qwen3.5-0.8B and one continuous
joint-training budget, with current-task and temporal losses active from the
first update. Its configured budget is eight hours on four GPUs, with per-rank
batches of 16 current observations and 32 directed temporal pairs and no gradient
accumulation. The resulting global batch sizes are 64 and 128, respectively;
these settings require capacity checks before a full run. Forecasting uses
source images and reports with signed elapsed time, without EHR inputs.
Evaluation of this implementation and matched native supervised comparisons
remain pending. Full six-task adaptation of that design, standard
MIMIC-CXR-VQA, and MS-CXR phrase grounding remain pending. All reported numbers
are point estimates; matched data and update counts do not imply matched GPU hours.
